\subsection{The touching upper bounds}
For an ordered facet word $\sigma=(\sigma_1,\ldots,\sigma_m)$, write
\[
 C_\sigma(a)=
 \begin{pmatrix}a_{\sigma_1}&\cdots&a_{\sigma_m}\\1&\cdots&1\end{pmatrix},
 \qquad e=(0,0,0,0,1)^T.
\]
The equations $C_\sigma(a)\beta=e$ impose closure and normalization.
For positive feasible weights let
\[
 Q_\sigma(a,\beta)=\sum_{j<i}\beta_i\beta_j
 \omega_0(a_{\sigma_j},a_{\sigma_i}),\qquad
 A_\sigma(a)=\frac1{2Q_\sigma(a,\beta(a))},\qquad
 U_\sigma(a)=\frac{A_\sigma(a)^2}{2V(a)}.
\]
When $Q_\sigma>0$, the variational capacity formula gives
$\operatorname{sys}(K(a))\leq U_\sigma(a)$.

Let $T$ be the tangent space at $a_0$ to the orbit generated by translations,
positive dilations and linear symplectic maps. The following lemma records
the finite calculation that will be used in the rest of the proof.
\begin{lemma}[Exact upper bounds and symmetry directions]\label{lem:hko-exact-bounds}
The space $T$ has dimension fifteen. On a neighborhood of $a_0$ there are
twenty-six smooth functions $U_j$ such that
\[
 \operatorname{sys}(K(a))\leq U_j(a),\qquad
 U_j(a_0)=\operatorname{sys}(K_0).
\]
Their derivative rows $r_j=DU_j(a_0)$ vanish on $T$, have rank twenty-five,
and satisfy
\[
 \sum_{j=1}^{26}\lambda_jr_j=0,\qquad
 \lambda_j>0,\qquad \sum_{j=1}^{26}\lambda_j=1.
\]
\end{lemma}
The construction and computer-algebra verification of this lemma occupy
the following calculations. Holding selected weights fixed and solving
the five closure and normalization equations for the others produces the
smooth sections: an invertible five-column minor gives smoothness, and
strict positivity at the base point persists nearby.

\subsection*{One of the twenty-six upper bounds}

We now work through one of the twenty-six upper bounds to illustrate the
construction and differentiation. The other twenty-five are constructed
in the same way. The exact computations for all twenty-six were verified
using a computer algebra system. Appendix~\ref{app:hko-cas} gives the complete
program with explanations of its checks.

% Corresponds to entry 2 (zero-based) of the computational witness.
Consider the seven-facet word and position sets
\[
 \sigma=(2,3,8,9,5,4,6),\quad I=(1,2,3,4,5),\quad J=(6,7).
\]
We solve closure and normalization for the weights at positions $I$ and
hold the weights at positions $J$ fixed. These are positions in the word,
not facet labels. Set
\[
 r=\frac{6909830056250513}{10^{17}},\quad
 c=\frac{t^2+5}{40},\quad z=\frac{5-t^2}{20}=\frac{s}{10},\quad
 k=\frac{3-t^2}{4}.
\]
Fix $\beta_6=r$ and $\beta_7=z$; the decimal-looking rational is an exact rational
choice, not a rounding assertion. The complete base vector is
\[
 \beta_0=(c+kr,\ c-r,\ c,\ c,\ z-kr,\ r,\ z)^T.
\]
At this word the closure matrix's selected minor is
\[
 M=C_\sigma(a_0)_I=
 \begin{pmatrix}
 -\alpha&-d&0&0&1\\b&0&0&0&-t\\0&0&0&-b&0\\
 0&0&d&\alpha&0\\1&1&1&1&1
 \end{pmatrix},\qquad \det M=10(1-t^2)>0.
\]
The full matrix is obtained by adjoining columns $(a_4,1)^T$ and $(a_6,1)^T$.
Writing $e=(0,0,0,0,1)^T$, define the nearby section by
\[
 \beta_J(a)=(r,z)^T,\qquad
 \beta_I(a)=C_\sigma(a)_I^{-1}\bigl(e-C_\sigma(a)_J(r,z)^T\bigr).
\]
Direct substitution gives $C_\sigma(a_0)\beta_0=e$ and $\beta_0>0$.
For example $0.726<t<0.727$ and $0.069<r<0.070$ suffice to check all
seven signs from the displayed formulas. Invertibility and positivity persist
in a neighborhood.

Use $\omega(x,y)=x_{q_1}y_{p_1}+x_{q_2}y_{p_2}-x_{p_1}y_{q_1}-x_{p_2}y_{q_2}$ and
\[
 Q_\sigma(a,\beta)=\sum_{j<i}\beta_i\beta_j\omega(a_{\sigma_j},a_{\sigma_i}).
\]
At the fixed base geometry, solving closure with arbitrary
$(\beta_6,\beta_7)=(u,v)$ gives the particularly simple identity
\[
 Q_\sigma=-4tv^2+(2t-2t^3/5)v.
\]
It is independent of $u$ and has its maximum at $v=z$; its value is $Q_0=t/5$.
Thus its Hessian on the two-dimensional closure affine space is singular
(in the $(u,v)$ coordinates it is $\operatorname{diag}(0,-8t)$).
The five-by-five closure minor is nevertheless invertible, so the displayed
feasible section continues smoothly and supplies an upper bound nearby.
\[
 A_0=\frac1{2Q_0}=\frac5{2t}=5t-\frac{t^3}{2},\qquad
 U_0=\frac{A_0^2}{2V_0}=\frac{11-t^2}{10}>1.
\]
The HKO capacity formula gives
$c_{\rm EHZ}(K(a_0))=2\cos(\pi/10)(1+\cos(\pi/5))=A_0$.
Thus the constructed upper bound agrees with the systolic ratio at the
base point.

\subsection*{A worked coordinate outside product-preserving perturbations}
We now choose a perturbation that varies the first momentum component of
the second $q$-facet's normal:
\[
 h_2=(0,0,1,0),\qquad h_i=0\quad(i\ne2).
\]
% This is flat coordinate 6 in the verifier's zero-based storage.
Put $w=c+kr=\beta_{0,1}$. Differentiated closure gives
\[
 D\beta_J[h]=0,\qquad D\beta_I[h]=-wM^{-1}(0,0,1,0,0)^T
 =w\begin{pmatrix}
 t^3/20-t/4\\t^3/20-t/4\\t^3/40-3t/8\\-3t^3/20+5t/4\\t^3/40-3t/8
 \end{pmatrix}.
\]
Inserting these values into the product rule for $Q$ gives the exact scalar
\[
 DQ[h]=\frac{3-t^2}{80}+
 \left(\frac{3t^2}{20}-\frac14\right)r
 -\frac{t^2+5}{4}r^2=:q_h.
\]
Since $a_2$ has zero momentum coordinates, $DV_0[h]=0$, while
\[
 DA[h]=-\frac{25}{2t^2}q_h,\qquad
 DU[h]=-\frac{125}{4t^3V_0}q_h.
\]
This coefficient of the upper-bound derivative is approximately
$-0.1801707325$. Changing only the two polygon factors keeps every
$q$-facet normal in the $q$-plane. This perturbation adds a $p_1$-component
to one of those normals, so the derivative describes a change of the upper
bound outside that product family.

\subsection{Computing the twenty-six derivatives}
\label{subsec:hko-local-maximum-computation-with-sagemath}
The table on page~\pageref{page:hko-exact-choices} gives the exact fixed values for every section,
including the worked entry. Except at entry 22, the minor positions are $I=(1,2,3,4,5)$ and
$J=(6,\ldots,m)$. Entry 22 instead has $I=(1,2,3,4,6)$ and $J=(5,7)$.
The table is ordered by the witness's zero-based certificate
index. For each table row set $\beta_J=u$ and compute
\[
 \beta_I=C_I^{-1}(e-C_Ju),\qquad
 (D\beta)_J[h]=0,\qquad
 (D\beta)_I[h]=-C_I^{-1}(DC[h])\beta.
\]
For each of the forty coordinate directions, compute
\begin{align*}
 DQ[h]&=\sum_{j<i}\bigl((D\beta_i[h])\beta_j+\beta_i(D\beta_j[h])\bigr)
       \omega(a_{\sigma_j},a_{\sigma_i})\\
 &\quad+\sum_{j<i}\beta_i\beta_j
 \bigl(\omega(h_{\sigma_j},a_{\sigma_i})+\omega(a_{\sigma_j},h_{\sigma_i})\bigr),\\
 R[h]&=-\frac{A_0}{2V_0Q_0^2}DQ[h]
       -\frac{A_0^2}{2V_0^2}DV_0[h].
\end{align*}
Together with the explicit table these formulas specify the entire
$26\times40$ matrix $R$ exactly.

The finite predicates are: every word has distinct valid labels;
$\det C_I\ne0$; $C\beta=e$; all $\beta_i>0$; $Q=t/5$;
the differentiated closure equations hold; the rows annihilate the fifteen
symmetry columns; $\operatorname{rank}R=25$; and the one-dimensional kernel
of $R^T$ has a vector whose entries have one strict sign. Normalize that
vector to sum to one, producing $\lambda>0$, $\lambda^TR=0$.
The symmetry columns are explicitly $h_i=-a_{i,j}a_i$ for $j=1,\ldots,4$,
$h_i=-a_i$, and $h_i=-X^Ta_i$ for the ten generators
\[
 X=\begin{pmatrix}A&B\\C&-A^T\end{pmatrix},\qquad B=B^T,\ C=C^T,
\]
using the four matrix units for $A$ and the three symmetric matrix units
for each of $B,C$. Their column rank is required to be fifteen.
These exact arithmetic checks establish the asserted independence, rank and
positive relation. The constructions above give smooth feasible sections
on a common neighborhood, and the HKO capacity formula gives the touching
identities. This proves Lemma~\ref{lem:hko-exact-bounds}.
